\section{Polarity, norm sums, and Hanner polytopes}\label{sec:convex}

All bodies below are symmetric about the origin. In addition to the
polar and volume product defined in Section~\ref{sec:introduction}, we use
the gauge and support function
\[
 p_A(x)=\inf\{t>0:x\in tA\},\qquad
 h_A(y)=\max_{x\in A}\langle x,y\rangle.
\]
Dual spaces are identified with Euclidean coordinate spaces, and
$|A|$ always denotes volume in the dimension of $A$. The notation $h_A$
also applies to any nonempty compact convex set. The gauge is a
norm with closed unit ball $A$. Convex separation gives the bipolar
identity $A^{\circ\circ}=A$, and hence
\begin{equation}\label{eq:gauge-support}
 p_A=h_{A^\circ},\qquad h_A=p_{A^\circ}.
\end{equation}
Indeed $x\in tA$ is equivalent to
$\langle x,y\rangle\le t$ for every $y\in A^\circ$.
For an invertible linear map $T$,
\begin{equation}\label{eq:linear-polarity}
 (TA)^\circ=T^{-\mathsf T}A^\circ,
 \qquad P(TA)=P(A),
\end{equation}
since the two volumes acquire reciprocal determinant factors.

For bodies $A\subset\R^k$ and $A'\subset\R^l$, define
\[
 A\oplus_\infty A'=A\times A',\qquad
 A\oplus_1 A'=\operatorname{conv}
       \bigl((A\times\{0\})\cup(\{0\}\times A')\bigr).
\]
These operations are taken in the product space. More generally,
complementary subspaces are identified with product coordinates by an
invertible linear map. Their gauges are
\begin{equation}\label{eq:sum-gauges}
 \begin{split}
 p_{A\oplus_\infty A'}(x,x')&=\max\{p_A(x),p_{A'}(x')\},\\
 p_{A\oplus_1 A'}(x,x')&=p_A(x)+p_{A'}(x').
 \end{split}
\end{equation}
For the second identity, a convex combination from the two coordinate
bodies has gauge sum at most one. Conversely, when the sum is at most
one, use these two gauges as the coefficients of the normalized
coordinate vectors and place any remaining weight at $0$. A zero
coordinate contributes no term.

\begin{lemma}[Polars and volumes of the two sums]\label{lem:sums}
For origin-symmetric convex bodies of dimensions $k,l\ge1$,
\begin{equation}\label{eq:sum-polars}
 (A\oplus_1 A')^\circ=A^\circ\oplus_\infty(A')^\circ,
 \qquad
 (A\oplus_\infty A')^\circ=A^\circ\oplus_1(A')^\circ.
\end{equation}
Moreover,
\begin{equation}\label{eq:sum-volumes}
 |A\oplus_\infty A'|=|A|\,|A'|,
 \qquad
 |A\oplus_1 A'|=\frac{k!\,l!}{(k+l)!}|A|\,|A'|.
\end{equation}
Consequently, for $s\in\{1,\infty\}$,
\begin{equation}\label{eq:sum-products}
 P(A\oplus_s A')=\frac{P(A)P(A')}{\binom{k+l}{k}}.
\end{equation}
\end{lemma}

\begin{proof}
Testing a functional on the two coordinate bodies gives the first
polar identity; taking polars again gives the second. The product
volume follows from Fubini. By \eqref{eq:sum-gauges}, the fiber of
$A\oplus_1 A'$ over $x\in A$ is $(1-p_A(x))A'$. Since
$|\{x:p_A(x)\le t\}|=t^k|A|$ for $0\le t\le1$, integration gives
\[
 |A\oplus_1 A'|
 =k|A|\,|A'|\int_0^1t^{k-1}(1-t)^l\,dt
 =\frac{k!\,l!}{(k+l)!}|A|\,|A'|.
\]
The last integral is evaluated by repeated integration by parts.
Combining the polar and volume identities proves
\eqref{eq:sum-products}.
\end{proof}

The Hanner recursion from Section~\ref{sec:introduction} uses precisely
these two operations, following Hanner~\cite{Hanner1956}. Their volume and polarity formulas give the
attaining family immediately.

\begin{lemma}[The Hanner volume product]\label{lem:hanner}
Every $d$-dimensional linear Hanner body $H$ satisfies
\[
 P(H)=\frac{4^d}{d!}.
\]
The polar of a linear Hanner body is again a linear Hanner body.
\end{lemma}

\begin{proof}
A centered interval $[-a,a]$, $a>0$, has polar $[-1/a,1/a]$ and
volume product $4$. Equation~\eqref{eq:sum-products} propagates
$4^k/k!$ and $4^l/l!$ to $4^{k+l}/(k+l)!$ under either operation.
Induction and linear invariance prove the first assertion. The polar
identities exchange the two operations and preserve the interval
base case. Equation~\eqref{eq:linear-polarity} then proves closure
under polarity also for linear images.
\end{proof}

In dimension one every origin-symmetric convex body is a centered
nondegenerate interval. Thus both the sharp inequality and the full
equality classification in that dimension follow directly from the
base case above.
